\section{固定量子VKの補完軌跡と再利用できる制限複体}

重み三のCGG共鳴 $(j,r)=(0,3),(1,1)$、$s\ne0$ に対し、
六項代数を $\Lambda$、$e=e_0+e_1+e_3+e_4$、$B=e\Lambda e$ とする。
固定した量子代数での結果は
\[
 HH^\bullet(B)=(1,1,35,39),\qquad
 R_B\simeq\mathbb C[[z_1,\ldots,z_{35}]],\qquad
 R_\Lambda\simeq\mathbb C[[\gamma]]/(\gamma^{j+1})
\]
である。適合的なhullでは制限を
\[
 \operatorname{Spf}R_\Lambda\hookrightarrow\operatorname{Spf}R_B
\]
と表示できる。$j=0$ では補完可能な一次方向は零であり、
$j=1$ では一次元だがその非零方向は二次で止まる。
可換基点の五十三次元の接空間同型とは、基点が異なる。

再利用すべき機構は、相対正規化複体の全射
\[
 C_E^\bullet(\Lambda_h)\longrightarrow C_{eE}^\bullet(B_h)
\]
とその核 $K_h^\bullet$ である。
可換基点での $HH^0,HH^1,HH^2$ の同型から $H^1(K_0)=H^2(K_0)=0$ が従う。
有限自由複体の階数特殊化によりgeneric fiberでも同じ消滅を得るため、
固定量子化での $HH^1$ 同型と $HH^2$ 単射が同時に得られる。
非零量子化係数の同型があれば、この結論を任意の $s\ne0$ へ移せる。

可換cornerの結合積スキームは次元 $53+598=651$ の滑らかな点を持つ。
量子化後もこの滑らかさが残り、軌道次元は
$620-1-3=616$ となるため、商のminiversal方向は三十五個である。
この計算は、障害群 $HH^3(B)$ が三十九次元残っていてもhullが滑らかになり得る
具体例として使える。

全 $\alpha$ に対して指定cornerを保つことは、単なる圏同値に頼らず確認できる。
\[
 a\star_{\alpha+\beta}b=T^{\beta\deg b}(a)\star_\alpha b,
 \qquad T=1+sN
\]
から、右加群の通常のEnd合成には $G_{ik}(a)=T^{-\beta i}(a)$ を使う。
経路順の積には $F_{ik}(a)=T^{\beta k}(a)$ を使う。
どちらも全冪等元を固定する。
Hochschild検算や既存quiver表示の規約を切り替える際、この向きの差を明示する。

補完像の差は、\emph{指定した外側と指定した埋め込み}を含む不変量である。
二つの $B$ 自体の導来非同値は、次節のGerstenhaber作用によって別に証明される。
今後、全ての最小補完を統一的に記述する関手を作る場合も、
この標準補完像と圏単体の不変量を区別する必要がある。
